\section{Related literature}

\subsection{CMUR construction, data quality and interpretive validity}

CMUR interpretation depends on the relationship between measured material flows and the ratio constructed from them. Eurostat's methodology defines that relationship, while work on material-flow accounting and uncertainty establishes the importance of source reliability and spatial comparability \parencite{Eurostat2018,FischerKowalski2011,Patricio2015}. The Ireland-focused assessment by \textcite{McCarthy2024} brings methodology, input quality and comparability directly into CMUR interpretation. This stream makes clear that a published indicator cannot be assessed independently of its measurement architecture. Temporal support adds a specific question: how much evidence that architecture supplies for the direction, timing and change in the pace of the published rate.

\subsection{CMUR trends, targets, nowcasts and forecasts}

Analyses of CMUR trajectories serve several temporal purposes. \textcite{Christis2023} assess trends and the prospects of doubling the rate, explicitly acknowledging biennial waste availability and linear interpolation between reporting years. \textcite{Nuss2024Nowcasting} address more timely estimates and outlooks. Their report documents Eurostat's estimation of missing odd years and explains why updates to underlying datasets can obstruct reproduction of a CMUR release unless matching source versions are used. These studies establish source-frequency awareness as part of CMUR analysis. Their focus is on trajectories, targets and indicator production; the present question is how source provenance conditions retrospective support for different temporal claims.

The closest temporal comparison is \textcite{Haider2026}. They restrict an EU-aggregate trend to waste-reporting years, examine the implications for fitted trend and nominal precision, and carry that subsample into forecasts. Their principal target-gap conclusion survives the treatment of estimated years. This is substantive partial overlap with the medium-term question here: the EU-level annual-versus-reporting-year comparison already exists. The residual empirical task is to assess common-rule direction and model-based precision across Member States, and to connect that comparison to claim-level annual change, turning and acceleration under explicitly bounded residual sensitivities. Aggregate forecast robustness alone does not establish that local temporal claims retain support.

\subsection{Temporal disaggregation and mixed-frequency statistics}

Official-statistics practice provides established methods for deriving higher-frequency series from lower-frequency observations and for reconciling series with temporal and accounting constraints \parencite{ESS2018,Quilis2018}. Mixed-frequency analysis also offers alternatives to forcing all inputs onto one observation schedule. For composite confidence indicators, \textcite{Mattera2023} use mixed-frequency regression to combine monthly information with quarterly economic activity. Their setting illustrates an established methodological response to unequal frequencies, rather than a CMUR-specific comparison.

These approaches concern the construction or estimation of a series from an information set. Our empirical question concerns the support for temporal interpretations of a series that is already disseminated. A complete annual sequence permits differences and turning points to be calculated; their support under source-constrained alternatives must be assessed separately. The distinction motivates retaining CMUR's released values at admissible anchors and evaluating how the temporal claims respond to explicit assumptions about intervening years.

\subsection{Sampling frequency and inferential fidelity}

Observation frequency can affect the fidelity of estimated temporal trends. \textcite{CortesGuzman2026} demonstrate this empirically by comparing annual river-invertebrate series with lower-resolution versions, finding changes in estimated trend direction and magnitude. Their evidence establishes a relevant inferential mechanism in environmental monitoring. CMUR poses a different information problem: a regularly disseminated composite incorporates a lower-frequency source. The ecological effect sizes therefore do not supply a quantitative correction for CMUR.

Together, these streams establish the indicator's source structure, its temporal uses, statistical methods for frequency mismatch and evidence that observation schedules can affect inference. They motivate a narrower empirical comparison across CMUR temporal claims: whether medium-term directional support and precision behave differently from annual-change, turning and acceleration support when admissible source provenance is respected. This is the gap addressed by the Member-State-wide assessment and its paired amplitude and threshold sensitivities.
